% Working prose started before P120 outcome inspection. Not a submitted manuscript.
\section{問題設定と関連研究}
LLMエージェントが外部APIを用いて複数段階の仕事を行うとき，過去の指示，行動，観測，推論を次の要求へ再入力する費用が蓄積する．必要な情報を残して入力を短縮するには，履歴の保持，状態の引継ぎ，ツール実行を制御するハーネスが必要になる．本研究の目的は，最終成果の性能を維持し，制御と失敗を含めた仕事全体の費用を減らすことである．本稿ではそのうち，一度だけの履歴要約と，その作成費用を取り上げる．

CodeAct\cite{codeact}型のエージェントでは，生成したPythonを実行し，変数を次の工程へ保持できる．一方，変数を保持しても，モデルへ毎回渡す会話履歴の費用は残る．MEM1\cite{mem1}は内部状態の更新と行動生成を結び付け，SKILL.state\cite{skillstate}は構造化状態を更新して履歴の増大を抑える．本稿の通常要約や同時出力を，これらと独立した新方式とは位置付けない．また，局所的なtoken短縮が実際の課金削減に直結しないことや，圧縮後の再取得と行動の変化は既に指摘されている\cite{costproxy}．

本稿は，状態保持型の同じエージェントに全保持，別呼出し要約，同じ応答での要約を実装し，固定予算で全仕事を比較する．要約の要求，実際の採用，次入力の短縮，後続要求全体の入力差，最終成功を区別する．保存した実行列で元履歴を復元する補助計数により，入力短縮と要約生成費用の関係も調べる．対象は既知４課題であり，一般的な品質保証や新しい圧縮原理の提案ではない．

\section{履歴と要約の制御}
\subsection{共通状態と全保持}
元の依頼，公開指示，過去のコード，返却された観測，推論を保存し，Pythonの状態を全工程で維持する．保存用の全履歴と，次のモデル要求へ渡す履歴を分離する．呼出しIDと返却IDの対応，長い観測の表示，予算通知，API実行口は共通とする．全保持では通常の履歴表示を継続する．

\subsection{別呼出し要約}
通常表示の入力が実tokenizerで16,384 token以上，かつ完結したやり取りが３回以上ある最初の境界で，一度だけ同じモデルへ要約を要求する．元のsystem指示と課題，直近２回のやり取りを残し，それより古い完結履歴を要約の対象とする．要約は完了事項，再利用可能な状態，未完了事項，不確実な点の４欄を持つ．各欄は最大32文字列，１文字列は最大2,048文字，compact JSON全体は8,192文字以内とする．構造の検査を通り，通常表示より次入力が厳密に短い場合だけ履歴へ採用する．元履歴とPython状態は削除しない．

\subsection{同じ応答での要約}
別呼出しによる履歴の再入力を避けるため，同じ起動境界で次の\texttt{run\_python}引数に要約欄を付ける．任意欄と，schema及び指示で必須とする欄を別群で比較する．特定toolの強制はせず，\texttt{tool\_choice=auto}を維持する．要約対象は古い完結履歴だけであり，今回のコードの実行結果を先取りして書かせない．今回のコードと実際の返却は，後続履歴に原文で残す．

通常の行動と要約は4,096出力tokenを共有する．有効な基礎引数のコードは元の参照IDと対応付けて一度だけ実行する．要約欄の欠落，不正，打切り，非採用のための再生成はしない．実行後に継続が必要な場合，同じ残予算通知で通常表示と要約表示を数え，短くなった場合だけ採用する．対応を確認した要約metadataは両表示から外すが，原記録には保持する．形式と長さの検査は，内容の真実性や課題への十分性を証明しない．

\section{評価手順と指標}
\subsection{全仕事の比較}
AppWorld\cite{appworld}の写真移動，電気代精算，返金，曲情報のCSV書出しの既知４課題を用いる．全保持対別呼出し，別呼出し対任意欄，別呼出し対必須欄の順に，各群８条件を実行する．各群を事前に固定し，その全結果を閉じてから次群を決めた．各条件は１回であり，同じ４課題の24実行を24個の独立課題とは数えない．未見の予約課題は使わない．

Mac Studio M3 Ultra・256 GiB，AppWorld 0.1.3.post1，MiniMax-M3 IQ4\_XS，llama.cpp b11669を用いる．context上限65,536，１スロット，温度は実行役1.0，要約役0.7，top-p 0.95，環境seed42，モデルseed57から役割別に増加させる．初回はcacheを無効として実測０を確認し，以後は通常のprefix cacheを用いる．課題別の条件順は順方向，逆方向，順方向，逆方向で固定する．

各条件は全役を合わせて80モデル要求，450,000論理token，1,800作業秒までとし，写真160，他80 APIまでとする．要求の上限時間600秒と出力上限4,096は，残予算にも従う．使用量不明の応答は実行せず，API効果と後処理が確認できた場合だけ次の固定条件へ進む．各最終状態を独立評価器で一度判定し，正解や評価結果を実行役へ返さない．失敗，欠測，未実行の枠も残す．

主指標は，全モデル要求についての論理入力と出力の和
\begin{equation}
 T=\sum_j (I_j+O_j)
\end{equation}
である．要約役，追加指示，schema，失敗と回復，cache済み入力も含む．API数，実時間，cache利用量は別に報告し，金額や計算量へ同一視しない．終端使用量の欠測は０や予約上限で補わない．採用判定は全件評価済み，使用量既知，監査一致，課題ごとの成功損失なし，共通成功１件以上，全４課題と共通成功課題の双方で総tokenが少ないことを要求する．必須欄の群では，実際の要約採用１件以上も求める．この開発判定は統計的な非劣性検定ではない．

\subsection{同じ保存行動列での入力差}
全24枠の全モデル要求を対象に，実際の入力と，採用された要約の代わりに元履歴を戻した入力を同じtokenizerで数える．後続のコード，応答，観測，残予算通知，出力上限，schemaを固定する．要約が未採用なら両入力は同じである．別要約の要求は両側に残し，同時出力のmetadataを再入力しない処理も両側へ適用する．

要求$j$での復元入力を$\widetilde I_j$とし，
\begin{equation}
 D=\sum_j(\widetilde I_j-I_j)
\end{equation}
を求める．別呼出し要約の実測入力と出力の和を$C$とし，$D-C$を同じ実行列上の入力差と作成費用の関係として示す．同時出力の要約部分だけの生成tokenは分離測定していないため，独自の$C$を推定しない．終端使用量が不明でも入力だけを計数できる場合があるが，元の総費用の欠測はそのまま残す．

この処理では新たなモデル生成，コード実行，API，課題評価を行わない．元の実測入力tokenとの一致を確認し，仮想入力のcontext上限と仕事予算の超過を明示する．履歴を戻せば実際の行動も変わり得るため，$D$や$D-C$は別方式の総費用，因果効果，品質維持の証拠ではない．直後の短縮量を残り呼出し数へ掛ける近似も行わない．
